\documentclass[a4paper,12pt]{article}
\usepackage[utf8]{inputenc}
\usepackage[T1]{fontenc}
\usepackage[french]{babel}
\usepackage{geometry}
\usepackage{listings}
\usepackage{xcolor}
\usepackage{enumitem}
\usepackage{titlesec}
\usepackage{fancyhdr}
\usepackage{amsmath}

\geometry{margin=2.5cm}

% Style pour le code
\lstset{
    backgroundcolor=\color{gray!10},
    basicstyle=\ttfamily\small,
    breaklines=true,
    frame=single,
    rulecolor=\color{gray!50},
    keywordstyle=\color{blue},
    commentstyle=\color{green!60!black},
    stringstyle=\color{orange},
}

% En-tête / pied de page
\pagestyle{fancy}
\fancyhf{}
\rhead{Admin Réseaux}
\lhead{\leftmark}
\cfoot{\thepage}

\title{\Huge\textbf{Admin Réseaux}\\[0.5em]\large Notes de cours}
\date{}

\begin{document}

\maketitle
\tableofcontents
\newpage

%-----------------------------------------------------------
\section{Rappels commandes Linux}
%-----------------------------------------------------------

\subsection{Structure d'une commande Bash}

Une commande Bash est constituée de :
\begin{itemize}
    \item un \textbf{programme}
    \item un \textbf{argument}
    \item une \textbf{option}
\end{itemize}

Exemple :
\begin{lstlisting}
ls /etc/apache2 -lah
ls -lah /etc/apache2   % equivalent, peuvent etre interchanges
\end{lstlisting}

\begin{itemize}
    \item \texttt{ls} : fait appel à un programme
    \item \texttt{/etc/apache2} : argument du programme
    \item \texttt{-lah} : une option
\end{itemize}

\textit{Si on met juste le programme, a-t-on un argument par défaut ?}\\
Oui, par défaut — pour \texttt{ls} c'est \texttt{.} (répertoire courant).

\subsection{Commandes de base}

\begin{description}[leftmargin=3cm, style=nextline]
    \item[\texttt{ls -lah}] Lister les fichiers (human-readable, all, droits)
    \item[\texttt{cd}] Changer de répertoire (\textit{change directory})\\
        Sans argument : revient à la racine utilisateur courant (\texttt{\$HOME})
    \item[\texttt{cat}] Concaténer / afficher le contenu d'un fichier
    \item[\texttt{mkdir}] Créer un répertoire (\textit{make directory})\\
        \texttt{-p} : crée les dossiers parents si nécessaire
    \item[\texttt{pwd}] Donne le chemin absolu courant
    \item[\texttt{passwd}] Modifier le mot de passe
    \item[\texttt{rm}] Supprimer (\textit{remove})\\
        \texttt{rm -Rf /*} \textbf{(à ne surtout pas faire !)} — \texttt{-R} récursif, \texttt{-f} force
    \item[\texttt{touch}] Créer un fichier vide
    \item[\texttt{nano}] Éditeur de texte
    \item[\texttt{cp <src> <dest>}] Copie de fichier ; \texttt{-R} pour récursif
    \item[\texttt{scp <src> <dest>}] Idem mais via SSH
    \item[\texttt{mv <src> <dest>}] Déplacer un fichier
    \item[\texttt{clear}] Vider le buffer / nettoyer le terminal
    \item[\texttt{usermod}] Ajouter un utilisateur dans un groupe
    \item[\texttt{useradd / adduser}] Créer un utilisateur
    \item[\texttt{groupadd / addgroup}] Créer un groupe
    \item[\texttt{sudo}] \textit{Substitute User Do} — par défaut l'utilisateur de substitution est \texttt{root}
    \item[\texttt{su}] \textit{Substitute User} : connexion temporaire à un autre utilisateur\\
        \texttt{su -} : root\\
        \texttt{su <nom\_utilisateur>}
\end{description}

%-----------------------------------------------------------
\section{Permissions Linux}
%-----------------------------------------------------------

\subsection{Lecture des permissions}

Exemple de sortie \texttt{ls -l} :
\begin{lstlisting}
-rw- rw- --x   root  root  script.sh
\end{lstlisting}

\begin{itemize}
    \item 1\textsuperscript{er} caractère : type d'élément
        \begin{itemize}
            \item \texttt{d} : directory
            \item \texttt{-} : fichier brut
            \item \texttt{l} : lien symbolique
        \end{itemize}
    \item Colonnes suivantes (par groupe de 3) : \textbf{utilisateur propriétaire} | \textbf{groupe propriétaire} | \textbf{autres utilisateurs}
\end{itemize}

\subsection{Permissions en octet}

\begin{center}
\begin{tabular}{|c|c|}
    \hline
    \textbf{Permission} & \textbf{Valeur} \\
    \hline
    r (read)    & 4 \\
    w (write)   & 2 \\
    x (execute) & 1 \\
    0 (rien)    & 0 \\
    \hline
\end{tabular}
\end{center}

\subsection{chmod — changer les permissions}

\begin{lstlisting}
chmod 0554 script.sh
       ^  groupe
  type d'elt
\end{lstlisting}

Option : \texttt{-R} pour appliquer récursivement.

\textbf{Bonne pratique :} ne jamais donner \texttt{w} et \texttt{x} simultanément — risque d'injection de code malveillant et d'exécution.

\begin{lstlisting}
chmod 0331 truc.exe   % DANGEREUX : ecriture + execution
chmod 0510 truc.exe   % ok
chmod +x   truc.exe   % A bannir (trop permissif)
\end{lstlisting}

Pour faire face : accorder soit \texttt{w} soit \texttt{x}, \textbf{jamais les deux}.

%-----------------------------------------------------------
\section{Scripts Bash}
%-----------------------------------------------------------

\subsection{Généralités}

Un script bash est une suite de commandes/instructions. Il permet de \textbf{limiter les erreurs humaines}.

\textbf{Règle :} on ne fait JAMAIS les MAS (Mises À jour planifiées) directement — il faut tester avant.

La \textbf{première ligne} doit toujours être le \textit{shebang} :
\begin{lstlisting}
#!/bin/bash
\end{lstlisting}

Un script peut être rendu exécutable avec \texttt{chmod} ou lancé via :
\begin{lstlisting}
bash script.sh
\end{lstlisting}

\subsection{Variables}

\begin{lstlisting}
var1="Hello"
\end{lstlisting}
En bash, les variables n'ont pas de type.

\subsection{PAM}

PAM dans l'environnement Linux : \textit{Pluggable Authentication Module} — plateforme d'auto-identification.

\subsection{Pipe et redirections}

\begin{lstlisting}
cat test.txt | tail -n k
\end{lstlisting}
Le \textbf{pipe} (\texttt{|}) permet d'envoyer le résultat d'une commande dans une autre (« tuyau »).\\
On peut enchaîner autant de pipes qu'on veut.

\medskip
\begin{tabular}{ll}
    \texttt{>}  & Redirige stdout vers un fichier \\
    \texttt{2>} & Redirige stderr vers un fichier \\
    \texttt{\&>} & Redirige stdout ET stderr vers un fichier \\
\end{tabular}

\subsection{Horodater un nom de fichier}

\begin{lstlisting}
touch "truc-$(date +%y%m%d%H%M).txt"
\end{lstlisting}

\subsection{Code de retour}

\begin{lstlisting}
echo $?
\end{lstlisting}
\begin{itemize}
    \item Retour \texttt{= 0} : commande bien exécutée, sans erreur
    \item Retour \texttt{$\neq$ 0} : erreur
\end{itemize}

Si code erreur = 1 et erreur = 2, le problème est une couche au-dessus.

\subsection{stdin, stdout, stderr}

\begin{itemize}
    \item \texttt{stdin}  : entrée standard
    \item \texttt{stdout} : sortie standard
    \item \texttt{stderr} : sortie d'erreur
\end{itemize}

\subsection{grep et expressions régulières}

\texttt{grep} affiche toute la ligne contenant le mot cherché.

Utilisation avec Regex (option \texttt{-E}) :
\begin{lstlisting}
grep -Eo '([0-9a-zA-Z]{1,3}\.){3}'
      ^  valeurs cherchees
      nb de car. par bloc
      separateur
      nb de fois
\end{lstlisting}
L'option \texttt{-o} n'affiche que ce qui est entre guillemets (la correspondance).\\
On peut ajouter une autre expression régulière.

\subsection{ping}

Sous Linux, \texttt{ping} ne s'arrête pas tout seul.

\begin{lstlisting}
ping ... -c n &    # -c n : nombre de requetes, & : arriere-plan
ping ... -c n | grep "suite de caracteres"
\end{lstlisting}

\subsection{Fonctions bash}

\begin{lstlisting}
nom() {
    echo "hello" $1   # $1 : 1er parametre
}

nom "truc en plus"    # affiche : hello truc en plus
\end{lstlisting}

\subsection{Conditions (if)}

\begin{lstlisting}
if [ x condition x ]; then   # espaces obligatoires entre crochets
    instruction
else / elif
    instruction
fi
\end{lstlisting}

Opérateurs de comparaison :
\begin{center}
\begin{tabular}{|c|l|}
    \hline
    \textbf{Opérateur} & \textbf{Signification} \\
    \hline
    \texttt{-eq} & equal \\
    \texttt{-lt} & lower than \\
    \texttt{-gt} & greater than \\
    \hline
\end{tabular}
\end{center}

Négation d'une condition :
\begin{lstlisting}
if [ ! -d $dossier ]   # -d : est-ce un dossier ?
\end{lstlisting}

\subsection{Boucle while}

\begin{lstlisting}
while [ condition ]; do
    instruction
done

sleep 2   # attendre 2 secondes
\end{lstlisting}

\subsection{Paramètres d'un script}

\begin{lstlisting}
# exemple.sh
echo "des trucs" $1 $2

# Appel :
./exemple.sh machin bidule
# Affiche : des trucs machin bidule
\end{lstlisting}

%-----------------------------------------------------------
\section{TP du 17/04 — Informations}
%-----------------------------------------------------------

\begin{itemize}
    \item Groupe et identifiant de l'utilisateur courant :
        \begin{itemize}
            \item \texttt{uid} pour l'utilisateur
            \item \texttt{gid} pour le groupe
        \end{itemize}
    \item Date sur 2 lignes différentes
    \item Nombre de fichiers dans un dossier : via pipe
    \item Horodater le fichier archivé
\end{itemize}

Commandes utiles :
\begin{lstlisting}
head -n -X        # supprime les X premieres lignes
ls -lah           # 40k = taille totale du repertoire courant (-l)
watch ls -l chemin  # voir les changements dans le dossier en direct
ls truc 2> log.log  # ecrit les codes d'erreur vers log.log
\end{lstlisting}

\textbf{Indication devoir :} $\approx 83$ lignes\\
Dossier de sortie : \texttt{/tmp} — fichier log\\
En cas d'échec : réessayage avec un nombre limité d'essais\\
\textbf{Rendu :} le fichier \texttt{.sh}

%-----------------------------------------------------------
\section{Sauvegarde}
%-----------------------------------------------------------

Outil Linux pour la sauvegarde :
\begin{itemize}
    \item \texttt{dd} : copie bit à bit (bytes ou bits)
    \item \textbf{Rsync} : sauvegarde fichiers
\end{itemize}

\begin{description}
    \item[Sauvegarde différentielle] Se base sur la sauvegarde complète
    \item[Sauvegarde incrémentielle] Se base sur la dernière sauvegarde
\end{description}

\subsection{Règle 3-2-1-0}

\begin{center}
\begin{tabular}{|c|l|}
    \hline
    \textbf{Chiffre} & \textbf{Signification} \\
    \hline
    3 & 3 sauvegardes \\
    2 & 2 supports différents \\
    1 & 1 hors site \\
    0 & 0 erreur de restauration \\
    \hline
\end{tabular}
\end{center}

La sauvegarde doit être \textbf{immuable} (ne peut pas être modifiée) et testée régulièrement.

\end{document}